\section*{第 \ref{ch:g9:arith} 章 --- 算术：因数与质数}

\begin{solution}{exo:g9:arith:1}
$36$ 的因数：$1, 2, 3, 4, 6, 9, 12, 18, 36$。
$45$ 的因数：$1, 3, 5, 9, 15, 45$。
$17$ 的因数：只有 $1$ 和 $17$（$17$ 是质数）。
\end{solution}

\begin{solution}{exo:g9:arith:2}
$2\,346$ 的末位数字是 $6$，能被 $2$ 整除，不能被 $5$ 整除。数字和
$2 + 3 + 4 + 6 = 15$，能被 $3$ 整除，不能被 $9$ 整除。末两位数字
为 $46$，$46 = 4 \times 11 + 2$，故不能被 $4$ 整除，因此 $2\,346$
不能被 $4$ 整除。
\end{solution}

\begin{solution}{exo:g9:arith:3}
$72 = 2^3 \times 3^2$; \quad
$150 = 2 \times 3 \times 5^2$; \quad
$210 = 2 \times 3 \times 5 \times 7$; \quad
$121 = 11^2$.
\end{solution}

\begin{solution}{exo:g9:arith:4}
判断 $101$ 是否为质数：只需检验质数 $p$，满足 $p^2 \leq 101$，即 $2, 3, 5, 7$，它们均
不能整除 $101$（为奇数，数字和为 $2$，末位不是 $0/5$，
$101 = 7 \times 14 + 3$）：$101$ 是质数。

$91 = 7 \times 13$，不是质数。

$143 = 11 \times 13$，不是质数。
\end{solution}

\begin{solution}{exo:g9:arith:5}
求 $48$ 与 $60$ 的公因数：48 的因数有 $1, 2, 3, 4, 6, 8,
12, 16, 24, 48$；$60$ 的因数有 $1, 2, 3, 4, 5, 6, 10, 12, 15, 20,
30, 60$；其中公因数为 $1, 2, 3, 4, 6, 12$，因此 $\gcd(48,60) = 12$。

分解质因数：$48 = 2^4 \times 3$ 与 $60 = 2^2 \times 3 \times 5$；
公共质因数取指数较小者：$2^2 \times 3 = 12$。
\end{solution}

\begin{solution}{exo:g9:arith:6}
求 $\gcd(255, 154)$：
\begin{align*}
255 &= 154 \times 1 + 101, \\
154 &= 101 \times 1 + 53, \\
101 &= 53 \times 1 + 48, \\
53 &= 48 \times 1 + 5, \\
48 &= 5 \times 9 + 3, \\
5 &= 3 \times 1 + 2, \\
3 &= 2 \times 1 + 1, \\
2 &= 1 \times 2 + 0 .
\end{align*}
最后一个非零余数为 1，因此 $\gcd(255, 154) = 1$（两数互质）。

求 $\gcd(1053, 325)$：
\begin{align*}
1053 &= 325 \times 3 + 78, \\
325 &= 78 \times 4 + 13, \\
78 &= 13 \times 6 + 0 .
\end{align*}
故 $\gcd(1053, 325) = 13$。
\end{solution}

\begin{solution}{exo:g9:arith:7}
由辗转相除法：$588 = 504 \times 1 + 84$；$504 = 84 \times 6 + 0$：
故 $\gcd(588, 504) = 84$。于是
\[
\frac{588}{504} = \frac{588 \div 84}{504 \div 84} = \frac{7}{6},
\]
即为既约分数。
\end{solution}

\begin{solution}{exo:g9:arith:8}
花束的数目必须同时整除 $84$ 和 $126$，最大可能值为
$\gcd(84, 126)$。分解质因数：$84 = 2^2 \times 3 \times
7$，$126 = 2 \times 3^2 \times 7$，故最大公因数为 $2 \times 3 \times 7 =
42$。她最多可扎 $42$ 束花，每束含
$\frac{84}{42} = 2$ 枝玫瑰、$\frac{126}{42} = 3$ 枝郁金香。
\end{solution}

\begin{solution}{exo:g9:arith:9}
需求最小公倍数。$24 = 2^3 \times 3$ 和
$36 = 2^2 \times 3^2$；各质因数取\emph{较大}
指数，得 $\lcm = 2^3 \times 3^2 = 72$。两班渡船下次同时发车
在 8:00 之后 $72$ 分钟，即 9:12。
\end{solution}

\begin{solution}{exo:g9:arith:10}
\emph{1.} 设 $d$ 为 $n$ 与 $n+1$ 的任一公因数，则它也整除它们的
差 $(n+1) - n = 1$，故 $d = 1$，即 $\gcd(n, n+1) = 1$。

\emph{2.} 分数为既约分数，当且仅当其分子与分母
互质，而 $n$ 与 $n + 1$ 互质，由第
1 小题可知。
\end{solution}

\begin{solution}{pb:g9:arith:1}
\textbf{1.} 将 $3$ 升壶装满，倒入 $5$ 升壶（水量
$0/3 \to 3$，在大壶中）。再将 $3$ 升壶装满，倒入
$5$ 升壶至满：大壶只能再装 $2$ 升，剩下
\[
3 - 2 = 1 \text{ L 在小壶中。}
\]
步骤：装满 $3$；倒 $3 \to 5$；装满 $3$；倒 $3 \to 5$。

\textbf{2.} 将 $5$ 升壶装满，倒入 $3$ 升壶（大壶中剩下 $2$ 升
）；倒空 $3$ 升壶，把 $2$ 升水倒入 $3$ 升壶；将 $5$ 升壶装满；
倒入 $3$ 升壶至满 --- 只需 $1$ 升，大壶中剩下
$\mathbf{4}$ L 在大壶中。共六步操作。

\textbf{3.} 所有这些量：$1$（第 1 题），$2$（第 2 题中两步操作之后），
$3$ 与 $5$（单独装满），$4$（
第 2 题），$6 = 3 + 3$（装满的小壶，再加上倒入
大壶的 $3$ 升），$7 = 5 + 2$，$8 = 5 + 3$（两壶都装满）。
从 $1$ 到 $8$ L 的每一个整数量，都可用 $5$ 与 $3$ 量出。

\textbf{4.} 开始时两壶的水量都是 $0$，是 $2$ 的倍数。
装满一壶使其水量变为 $6$ 或 $4$，都是偶数。倒空一壶使其水量变为
$0$，也是偶数。倒水是在水量均为偶数的两壶之间转移水，
所倒的水量是两个偶数之差
（剩余空间，或可用水量），于是各壶水量始终为偶数。
奇数目标如 $1$ L 是无法达到的。

\textbf{5.} $\gcd(6, 4) = 2$：只能得到偶数量，这一点已由第 4 题
证实。$\gcd(5, 3) = 1$：该法则允许每个整数量，
第 3 题已全部实现。最大公因数恰好是
两壶的计量单位。

\textbf{6.} $91 = 1 \times 65 + 26$; $65 = 2 \times 26 + 13$;
$26 = 2 \times 13 + 0$：$\gcd(91, 65) = 13$。再看
$2\,026 = 44 \times 46 + 2$; $46 = 23 \times 2 + 0$:
$\gcd(2\,026, 46) = 2$.

\textbf{7.} 反复将 $5$ 升倒入 $13$ 升壶：两次装满后
大壶中有 $10$ 升；第三次装满只能装入
$3$，剩下 $5 - 3 = 2$ 留在小壶中 --- 即
$13$ 被 $5$ 除的余数是 $3$，而数量 $3$（空位）与 $2$
（剩余量）恰好是欧几里得的数（$13 = 2 \times 5 + 3$，
$5 = 1 \times 3 + 2$）。接着出现 $3 - 2 = 1$，这正是
算法的下一个余数。喷泉用水完成了欧几里得的
除法运算。

\textbf{8.} $\gcd(13, 5) = 1$，因此由第 5 题的法则，
每个整数量都可量出 --- 而第 7 题的倒水过程
确实得到了 $1$ L。是的。

\textbf{9.} $n$ 与 $2n + 1$ 的任一公因数整除
$2n + 1 - 2 \times n = 1$，故它必为 $1$。因此
$\gcd(n, 2n+1) = 1$ 总成立。

\textbf{10.} A 车：7:12，7:24，7:36，7:48，8:00\,\dots；B 车：
7:18，7:36，7:54\,\dots\ 首次同时发车：7:36，经过
$36$ 分钟 --- 也就是 $12$ 与 $18$ 的第一个公倍数。
规律：$36 \times \gcd(12, 18) = 36 \times 6 = 216 = 12 \times
18$。对 $5$ 与 $3$，第一个公倍数为 $15$，并有
$15 \times \gcd(5,3) = 15 \times 1 = 15 = 5 \times 3$.

\textbf{11.} 周期为 $17$ 年时，下一次重合
即 $17$ 与 $4$ 的第一个公倍数；因为
$\gcd(17, 4) = 1$，故为 $17 \times 4 = 68$ 年 --- 蝉
每四次出现才与高峰相遇一次。若周期为
$16$ 年，则 $16$ 是 $4$ 的倍数，故 \emph{每一次}
出现都恰遇高峰。质数长度的周期与任何更短的捕食者周期
都没有公因子，使重合的间隔尽可能拉长：
算术即伪装。

\textbf{12.} $2$ 的指数有四种选择，$3$ 的有三种，
$5$ 的有两种：$4 \times 3 \times 2 = 24$ 个因数。

\textbf{13.} 若 $m = 2^{a} \times 3^{b} \times \cdots$，则
$m^2 = 2^{2a} \times 3^{2b} \times \cdots$：每个指数都
加倍，因而为偶数。在 $360 = 2^3 \times 3^2 \times 5$ 中，
$2$ 与 $5$ 的指数均为奇数，所以 $360$ 不是完全
平方数。

\textbf{14.} 一个因数与自身配对，当且仅当
$d = \frac nd$，即 $n = d^2$：只有平方数才有这样的
中间因数。对于其他所有 $n$，因数可分成
若干对，个数为偶数。因此：因数个数为奇数
$\Leftrightarrow$ 完全平方数。验证：$36$ 的因数有
$1, 2, 3, 4, 6, 9, 12, 18, 36$ --- 共 9 个，为奇数，且
$36 = 6^2$；而 $360$ 有 $24$ 个因数（第 12 题），个数为偶数，且
不是完全平方数（第 13 题）。

\textbf{15.} 储物柜 $n$ 被每位学生 $k$ 各拨动一次，
其编号整除 $n$：总拨动次数等于 $n$ 的
因数个数。储物柜最终\emph{打开}，当且仅当被拨动的
次数为奇数 --- 由第 14 题，恰在 $n$ 为
完全平方数时。打开的储物柜为 $1, 4, 9, 16, 25, 36, 49, 64,
81, 100$ --- 共 10 个。
\end{solution}
